\documentclass[10pt]{mypackage}

\usepackage[uprightgreeks,noDcommand]{kpfonts}
\renewcommand{\coloneq}{\coloneqq}
\renewcommand{\eqcolon}{\eqqcolon}
%\usepackage{dsfont}
%\renewcommand{\mathbb}{\mathds}
%\usepackage{homework}
\usepackage{notes}

\usepackage[ backend=bibtex, style = alphabetic, sorting=ynt ]{biblatex}
\addbibresource{../all_references.bib}

\usepackage{parskip}
\frenchspacing

\fancyhf{}
\fancyhead[R]{Avinash Iyer}
\fancyhead[L]{Gromov's Theorem on Groups of Polynomial Growth}
\fancyfoot[C]{\thepage}

\setcounter{secnumdepth}{0}

\begin{document}
\RaggedRight
This is going to be a broad overview of the ideas and techniques used to prove the following much celebrated theorem.
\begin{theorem}[Gromov's Theorem on Groups of Polynomial Growth]
  If a finitely generated group $\Gamma$ has polynomial growth, then $\Gamma$ contains a nilpotent subgroup of finite index.
\end{theorem}
All of this work is discussed in Gromov's paper, and I am only writing this up for personal use (and to present to the University of Virginia Graduate Students Seminar).
\section{Notions from Riemannian Geometry}%
\begin{definition}
  Let $f\colon X\rightarrow Y$ be a map. We call such a map \textit{globally expanding} if, for any two points $x_1\neq x_2$ in $X$, one has $d\left( f\left( x_1 \right),f\left( x_2 \right) \right) > d\left( x_1,x_2 \right)$.

  We say $f$ is expanding if each point $x\in X$ has a neighborhood $U$ such that $f|_{U}$ is globally expanding.
\end{definition}
Consider a compact connected Riemannian manifold, and let $f\colon X\to X$ be an expanding map. Then, $X$ has no boundary. There is then a map $ \widetilde{f}\colon \widetilde{X}\rightarrow \widetilde{X} $ induced on the universal cover. Since any globally expanding map is a covering map on itself, it follows that the map $ \widetilde{f} $ is a homeomorphism, and the inverse map $ \widetilde{f}^{-1}\colon \widetilde{X}\rightarrow \widetilde{X} $ is a contraction.

Furthermore, for any $\delta > 0$, there is some $\ve > 0$ such that for any $y_1,y_2\in \widetilde{X}$ with $d_{\widetilde{X}}\left( y_1,y_2 \right) \geq \delta$, one has
\begin{align*}
  d_{\widetilde{X}} \left( \widetilde{f}^{-1}\left( y_1 \right),\widetilde{f}^{-1}\left( y_2 \right) \right) &\leq \left( 1-\ve \right)d_{\widetilde{X}} \left( y_1,y_2 \right),
\end{align*}
which follows by the definition of globally expanding. Therefore, $ \widetilde{f}^{-1} $ has a unique fixed point by the contraction mapping theorem, so projecting back down yields that $f\colon X\rightarrow X$ has a fixed point as well.

Two theorems of Shub and Franks are fundamental to the proof. To understand them, though, we require understanding a few results.
\section{Growth Rates and some Established Results}%
\begin{definition}
  Suppose $\Gamma$ has polynomial growth. We will let $\operatorname{growth}\left( \Gamma \right)$ be the smallest natural number $d$ such that there exists $C$ such that $\left\vert B\left( r \right) \right\vert\leq Cr^{d}$ for each $r\geq 1$. Here, $B\left( r \right)$ denotes the closed unit ball of radius $r$ centered at the identity.
\end{definition}
Suppose $\Gamma'\leq \Gamma$ is a subgroup of a finitely generated group acting via translation. Let $X = \Gamma/\Gamma'$ be the quotient space, and let $f\colon \Gamma\to X$ the natural projection. Let $\dist\left( x,y \right) = \inf_{\alpha,\beta}d\left( \alpha,\beta \right)$ for any $\alpha\in f^{-1}\left( x \right)$ and $\beta\in f^{-1}\left( y \right)$. Since $\Gamma'$ acts via isometries on $\Gamma$, it follows that $\dist\left( x,y \right)$ is a metric on $X$.

Observe that for any two points $x,y\in X$ with $\dist\left( x,y \right) = p$ for some integer $p$, there exist points $x = x_0,x_1,\dots,x_p=y$ such that $\dist\left( x_i,x_{i-1} \right) = 1$. Importantly, this means that if $X$ is infinite, then $\left\vert B_X(x,r) \right\vert\geq r+1$ for all $x\in X$ and all $r\geq 0$.
\begin{lemma}[Splitting Lemma]
  If $\Gamma'\subseteq \Gamma$ is a \textit{finitely generated} subgroup of \textit{infinite} index. Then, $\operatorname{growth}\left( \Gamma' \right)\leq \operatorname{growth}\left( \Gamma \right)-1$.
\end{lemma}
\begin{proof}
  Since $\Gamma/\Gamma'$ is connected, it follows that $B(r)$ for each $r\in \Gamma$ contains at least $r+1$ elements $\alpha_0,\dots,\alpha_r\in \Gamma$ such that $f\left( \alpha_i \right)\neq f\left( \alpha_j \right)$ whenever $i\neq j$, where $f\colon \Gamma\to \Gamma/\Gamma'$ is the quotient map. The intersections $B' = B(r)\cap \Gamma'$ and the translates $B'\alpha_i$ are disjoint, and they are contained in $B\left( 2r \right)$, so that $\left\vert B\left( 2r \right) \right\vert\geq \left( r+1 \right)\left\vert B' \right\vert$.
\end{proof}
\section{Linear Representations of Certain Finitely Generated Groups}%
\begin{theorem}[Jordan's Theorem]
  For any Lie group $L$ with finitely many components, there is a number $q$ such that every finite subgroup in $L$ contains an abelian subgroup of index at most $q$.
\end{theorem}
\begin{theorem}[Tits Alternative]
  Let $L$ be a Lie group with finitely many components, and let $G\subseteq L$ be an arbitrary finitely generated subgroup. Then, either
  \begin{itemize}
    \item $G$ has a free group of rank $2$ (hence is of exponential growth);
    \item $G$ is virtually nilpotent (hence $G$ has exponential growth unless it is virtually nilpotent).
  \end{itemize}
\end{theorem}
\begin{theorem}[Milnor--Wolf Theorem]
  If $\Gamma$ is a virtually solvable group, then either $\Gamma$ is of exponential growth or $\Gamma$ is virtually nilpotent.
\end{theorem}
\begin{remark}
  Given a nilpotent group $G$ with lower central series
  \begin{align*}
    G = \gamma_1(G) \geq \gamma_2(G)\geq\cdots\geq \gamma_c(G)\geq \gamma_{c+1}(G) = \set{e_G},
  \end{align*}
  the Bass--Guivarch formula gives that $G$ grows polynomially with degree given by
  \begin{align*}
    d &= \sum_{i=1}^{c}i\operatorname{rk}\left( \gamma_i(G)/\gamma_{i+1}(G) \right),
  \end{align*}
  where $\operatorname{rk}$ denotes the rank of the abelian groups $\gamma_i(G)/\gamma_{i+1}(G)$.
\end{remark}
\begin{lemma}[Lemma (a)]
  Let $L$ be a Lie group with finitely many components, and suppose that for each $p=1,2,\dots$, there is a homomorphism $G\to L$ such that its image is finite and has at least $p$ elements. Then, $G$ contains a subgroup $G'\subseteq G$ of finite index such that $\left[ G',G' \right]\subseteq G'$ has infinite index.
\end{lemma}
\begin{proof}
  We start by observing that if $G$ is a finitely generated group with finite generating set $S$, then given any index $k$, $G$ has finitely many subgroups of index $k$. This follows from the fact that if $H$ is any such subgroup, then the action of $G$ on $G/H$ yields a homomorphism $G\to S_k$ (and so $H$ is determined by the given homomorphism), of which there are at most $\left( k! \right)^{\left\vert S \right\vert}$ such homomorphisms.

  Therefore, there are finitely many subgroups of index at most $q$, and since the finite intersection of finite index subgroups is finite index, it follows that if we let
  \begin{align*}
    G' &\coloneq \bigcap\set{H\leq G : \left[ G:H \right]\leq q},
  \end{align*}
  then $G'$ is of finite index.

  We will show that $\left[ G',G' \right]$ is of infinite index by showing that $G'$ admits a surjection onto abelian subgroups of arbitrary size. If $\phi\colon G\to L$ has image $\left\vert F_p \right\vert\geq p$, then $G'\subseteq \phi^{-1}\left( A_p \right) $, where $A_p\leq F_p$ is abelian and has index $q$. Since
  \begin{align*}
    \left[ F_p:\phi\left( G' \right) \right] &\leq \left[ G:G' \right],
  \end{align*}
  it follows that
  \begin{align*}
    \left\vert \phi\left( G' \right) \right\vert &\geq \frac{p}{\left[ G:G' \right]}.
  \end{align*}
  Since $p$ is arbitrary, it follows that $\left[ G',G' \right]\subseteq G'$ has infinite index.
\end{proof}
\begin{lemma}[Lemma (b)]
  Suppose $\Gamma$ is a finitely generated group of polynomial growth. Then, $\left[ \Gamma,\Gamma \right]$ is finitely generated.
\end{lemma}
\begin{proof}
  It suffices to show that the kernel $\Delta\subseteq \Gamma$ of any homomorphism $g\colon \Gamma\to\Z$ is finitely generated. 

  Let $\gamma_0,\gamma_1,\dots,\gamma_k\in \Gamma$ be a system of generators such that $g\left( \gamma_0 \right) = 1\in \Z$ and $\gamma_i\in \Delta$ for each $i=1,\dots,k$. Let
  \begin{align*}
    \Delta_m &= \left\langle \set{\gamma_0^{j}\gamma_i\gamma_0^{-j} : 1\leq i \leq k, -m\leq j \leq m} \right\rangle.
  \end{align*}
  Then,
  \begin{align*}
    \Delta &= \bigcup_{m=0}^{\infty} \Delta_m.
  \end{align*}
  If the sequence of $\Delta_m$ stabilizes, then it equals $\Delta$ and we are done. Else, there are $\alpha_m\in \Delta$ such that each $\alpha_m$ is of the form $\alpha_m = \gamma_0^{m}\gamma_i\gamma_0^{-m}$ or $\alpha_m = \gamma_0^{-m}\gamma_i\gamma_0^{m}$ satisfying $\alpha_m\notin \left\langle \alpha_0,\alpha_1,\dots,\alpha_{m-1} \right\rangle$.

  Were such a family to exist, our goal is to show that it implies the existence of an exponentially growing family. Let $\beta\left( \epsilon_0,\dots,\epsilon_m \right) = \alpha_0^{\epsilon_0}\cdots\alpha_m^{\epsilon_m}$, where $\epsilon_i \in\set{0,1}$ for each $i$. Since equality between any two such elements $\beta$ is defined by equality of each $\epsilon_i$, it follows that there are $2^{m+1}$ such $\beta$.

  Notice then that $\ell\left( \beta \right)\leq \sum_{i=1}^{m}\ell\left( \alpha_m \right)\leq \left( m+1 \right)\left( 2m+1 \right)$. Thus, $\left\vert B\left( \left( m+1 \right)\left( 2m+1 \right) \right) \right\vert\leq 2^{m+1}$ as the ball of radius $\left( m+1 \right)\left( 2m+1 \right)$ has all the $\beta$. Yet, this contradicts polynomial growth.
\end{proof}
The most important applications of these are to prove the following lemma. 
\begin{lemma}[Lemma 1]
  Let $\Gamma$ be a finitely generated group of polynomial growth, and let $L$ be a Lie group with finitely many connected components. Suppose that for each finitely generated infinite group $\Gamma'\subseteq \Gamma$, there is a subgroup $\Delta\subseteq \Gamma'$ of finite index such that, for all $p=1,2,\dots$, there is a homomorphism $\Delta\to L$ such that its image contains at least $p$ elements.

  Then, $\Gamma$ is virtually nilpotent.
\end{lemma}
\begin{proof}
  By the splitting lemma, we inductively assume that all finitely generated subgroups of $\Gamma$ with infinite index are virtually nilpotent.

  Let $\Delta\subseteq \Gamma$ be a subgroup of finite index with the required homomorphisms into $L$. If these homomorphisms are all finite, Lemma (a) yields a subgroup $\Delta'\subseteq \Delta$ with finite index such that $\left[ \Delta',\Delta' \right]\subseteq \Delta$ has infinite index. Induction yields the final result in this case.

  If there is a homomorphism $\Delta\to L$ with infinite image, then the Tits theorem applied to this image yields a finite index subgroup $\Delta'\subseteq \Delta$ such that $\left[ \Delta',\Delta' \right]\subseteq \Delta$ is infinite index as well. According to Lemma (b), the commutator $\left[ \Delta',\Delta' \right]$ is finitely generated, so by the induction hypothesis, it is virtually nilpotent. Therefore, $\Gamma$ is virtually solvable, so by the Milnor--Wolf Theorem, it follows that $\Gamma$ is virtually nilpotent.
\end{proof}
\section{Constructing the Asymptotic Cone}%
In Gromov's original paper, he constructed a special metric space $Y$ as a subsequential limit of the sequence of discrete metric spaces $\left( \Gamma,\left( 1/n \right)d_S \right)$, where $S$ is a generating set for $S$. The metric space $Y$ has the following properties:
\begin{enumerate}[(i)]
  \item $Y$ is homogeneous;
  \item $Y$ is connected and locally connected;
  \item $Y$ is complete;
  \item $Y$ is locally compact and of finite topological dimension.
\end{enumerate}
The deep result of Gleason--Montgomery--Zippin gives that $\operatorname{Isom}\left( Y \right)$ is then a connected Lie group. There is then a finite-index subgroup $\Gamma'\subseteq \Gamma$ that admits a homomorphism $\Gamma'\to \operatorname{Isom}\left( Y \right)$, and we use Lemma 1 to then obtain the desired result for $\Gamma$.
\begin{definition}
  Given a metric space $Y$ with a distinguished point $e$, $\operatorname{Isom}\left( Y \right)$ is topologized by defining the basic closed neighborhoods
  \begin{align*}
    U_{k,\ve} &= \set{\sigma\in \operatorname{Isom}\left( Y \right) : k\in \Z_{>0}\text{ and }\ve > 0\text{ with }d\left( \sigma(y),y \right) \leq \ve\text{ for all $y$ with }d\left( y,e \right)\leq k}
  \end{align*}
  at the identity.
\end{definition}
The way we will construct this is through the asymptotic cone.
\subsection{The Nonstandard Reals}%
Let $\R$ denote the real numbers, and let $ \mathcal{U} $ be an ultrafilter on $\N$. On the set $\R^{\N}$, we may define an equivalence relation by saying two sequences $\left( u_n \right)_n$ and $\left( v_n \right)_n$ are equivalent if
\begin{align*}
  \set{n\in \N : u_n = v_n}&\in \mathcal{U}.
\end{align*}
We denote by $\R^{\mathcal{U}}$, which will be written as $\R^{\ast}$, the ultrapower of $\R$. It c an be shown that the operations that turn $\R$ into an ordered field, such as addition, subtraction, multiplication, division, and the ordering on $\R$, are elementary expressions, so by the Løs theorem, they admit extensions to $\R^{\ast}$ that turn $\R^{\ast}$ into an ordered field. We will identify $\R\subseteq \R^{\ast}$ by taking $s\mapsto \left[ s \right]$, where $\left[ s \right]$ is the equivalence class of the constant sequence with value $s$.

In particular, we define
\begin{align*}
  \left[ \left( u_n \right)_n \right]\pm \left[ \left( u_n \right)_n \right] &= \left[ \left( u_n \pm v_n \right)_n \right]\\
  \left[ \left( u_n \right)_n \right] \left[ \left( v_n \right)_n \right] &= \left[ \left( u_nv_n \right)_n \right]\\
  \left[ \left( u_n \right)_n \right]/\left[ \left( v_n \right)_n \right] &= \begin{cases}
    u_n/v_n & v_n\neq 0\\
    0 & v_n = 0,
  \end{cases}
\end{align*}
and we say that $\left[ \left( u_n \right)_n \right]\leq \left[ \left( v_n \right)_n \right]$ if and only if $\set{n : u_n\leq v_n}\in \mathcal{U}$.

We note that the nonstandard construction does not preserve the Archimedean property. In particular, we may define an element $\left[ u \right]$ to be \textit{finite} if there exists $s\in \R$ with $s > 0$ such that $\left[ -s \right]\leq \left[ u \right]\leq \left[ s \right]$, and we say the element is \textit{infinite} if it is not finite. We say the element is \textit{infinitesimal} if for all $s\in \R$ with $s > 0$, we have $\left[ -s \right]\leq \left[ u \right]\leq \left[ s \right]$. If $\R^{\operatorname{fin}}$ denotes the finite elements of $\R^{\ast}$ and $\R^{0}$ denotes the infinitesimal elements, then there is an isomorphism $\rho\colon \R\to \R^{\operatorname{fin}}/\R^{0}$; if $h\colon \R\to \R^{\operatorname{fin}}$ is the natural inclusion map, then the homomorphism $\operatorname{st} = \rho^{-1}\circ h$ from $\R^{\operatorname{fin}}$ to $\R^{0}$ is called the \textit{standard part} map. 
\subsection{Constructing the Space $Y$}%
\begin{definition}
  Let $\omega$ be an ultrafilter on $\N$, and let $\Gamma$ be a finitely generated group. We will let $\Gamma^{(M)}$ denote the space of sequences 
\end{definition}
\nocite{gromov_polynomial_growth_paper,van_den_dries_wilkie_paper,topics_in_groups_and_geometry}
\printbibliography
\end{document}
